\section{An extension to convex Bayes-risk generators}
\label{sec:general-losses}

The logarithmic example has a broader counterpart for excess prediction
loss. This extension concerns the Jensen gap of a convex generator; it
does not assert the total-risk conclusion proved for logarithmic loss.
In particular, additive affine terms do not affect its distortion.

Let $\phi:[0,1]\to\mathbb R$ be continuous and strictly convex, and
differentiable on $(0,1)$. For a finite source with masses $p_x>0$ and
binary-label posteriors $\eta_x\in(0,1)$, define
\begin{equation}
 D_\phi(\mathcal P)
 =\sum_{C\in\mathcal P}\left\{
       \sum_{x\in C}p_x\phi(\eta_x)
       -p_C\phi(\bar\eta_C)\right\},
 \qquad
 p_C=\sum_{x\in C}p_x,\quad
 \bar\eta_C=p_C^{-1}\sum_{x\in C}p_x\eta_x.
 \label{eq:gen-distortion}
\end{equation}
This is the excess Bayes risk when the concave Bayes-risk function is
$-\phi$, up to an affine term, as in the convex representation of proper
scoring rules \cite[Theorem 1]{gneiting2007}. The identity follows by subtracting the
Bayes risk with the full source from that with the partition label.
Equivalently, putting
\begin{equation}
 d_\phi(z\Vert x)
 =\phi(z)-\phi(x)-\phi'(x)(z-x),
 \label{eq:gen-bregman}
\end{equation}
the distortion in \eqref{eq:gen-distortion} equals
$\sum_C\sum_{x\in C}p_xd_\phi(\eta_x\Vert\bar\eta_C)$, because the
linear terms sum to zero within every cell.

Write $D_{\phi,k}^*$ for the minimum of
\eqref{eq:gen-distortion} over partitions with $k$ nonempty cells.
For a four-state source, let
\begin{equation}
 \rho_\phi
 =\min_{\mathcal H}
   \max_{k\in\{2,3\}}
       \frac{D_\phi(\mathcal H_k)}{D_{\phi,k}^*},
 \label{eq:gen-price}
\end{equation}
where $\mathcal H$ ranges over complete binary-merger hierarchies.
The constructions below have four distinct posteriors, so both
denominators are strictly positive.

\begin{theorem}[Divergent endpoint slope]
\label{thm:gen-endpoint}
Suppose, in addition, that $\phi(x)=\phi(1-x)$ and
$\phi'(x)\to-\infty$ as $x\downarrow0$. Put
\[
 C=\phi(0)-\phi(1/2)>0,\qquad
 s_\varepsilon=-\phi'(\varepsilon),\qquad
 \delta_\varepsilon=\sqrt{C/s_\varepsilon}.
\]
There is a choice of positive $w_\varepsilon\le\varepsilon^2$ such
that the four-state source with masses
\[
 (p_A,p_C,p_D,p_B)=(c,w,w,c),\qquad c=(1-2w)/2,
\]
and posteriors
\[
 (\eta_A,\eta_C,\eta_D,\eta_B)
  =(\varepsilon,\delta,1-\delta,1-\varepsilon),
 \qquad w=w_\varepsilon,\quad\delta=\delta_\varepsilon,
\]
satisfies
\begin{equation}
 \rho_\phi
  \sim \frac{\sqrt{s_\varepsilon}}{2\sqrt C}
  \longrightarrow\infty
  \qquad(\varepsilon\downarrow0).
 \label{eq:gen-rate}
\end{equation}
All source and conditional label probabilities are positive for
sufficiently small $\varepsilon$.
\end{theorem}

\begin{proof}
We first justify a choice of $w$ that makes the required expansions
valid even when $\phi$ has no uniform second-derivative bound near
zero. For positive weights $\alpha,\beta$, define the weighted Jensen
gap
\[
 G(\alpha,\beta;x,z)
 =\alpha\phi(x)+\beta\phi(z)
  -(\alpha+\beta)\phi\left(
       \frac{\alpha x+\beta z}{\alpha+\beta}\right).
\]
If $m=(\alpha x+\beta z)/(\alpha+\beta)$, direct cancellation gives
\begin{equation}
 G(\alpha,\beta;x,z)
  =\beta d_\phi(z\Vert x)
       -(\alpha+\beta)d_\phi(m\Vert x).
 \label{eq:gen-small-weight}
\end{equation}
Fix an interior $\varepsilon$. Differentiability at $\varepsilon$
means
$d_\phi(\varepsilon+h\Vert\varepsilon)=o(|h|)$ as $h\to0$.
For $\alpha=c=(1-2w)/2$, $\beta=kw$, $k\in\{1,2\}$, and
$z\in[0,1]$, one has $|m-\varepsilon|\le4w$ for sufficiently
small $w$. Thus the remainder in \eqref{eq:gen-small-weight} is
$o(w)$ uniformly over these $k$ and $z$, with $\varepsilon$ fixed.
Since $s_\varepsilon\to\infty$, we may choose, separately for each
sufficiently small $\varepsilon$, a positive
$w=w_\varepsilon\le\varepsilon^2$ for which
\begin{equation}
 0\le
 kwd_\phi(z\Vert\varepsilon)
  -G(c,kw;\varepsilon,z)
 \le w/s_\varepsilon
 \quad
 (k\in\{1,2\},\ z\in[0,1]).
 \label{eq:gen-remainder-choice}
\end{equation}
The uniform statement follows directly from the definition of
differentiability on the shrinking interval
$[\varepsilon-4w,\varepsilon+4w]$; $w$ can be reduced further so
this interval is interior. This is an existence choice of the rare
mass, rather than an assumption of uniform differentiability as
$\varepsilon$ changes.

There are two useful preliminary limits. The supporting-line
inequality at $\varepsilon$ gives
\[
 0\le\varepsilon s_\varepsilon
 \le2\{\phi(\varepsilon/2)-\phi(\varepsilon)\}
 \longrightarrow0.
\]
It follows that $ws_\varepsilon\to0$ and
$\varepsilon/\delta\to0$, while
$\delta\to0$ and $\delta s_\varepsilon=\sqrt{Cs_\varepsilon}
\to\infty$. Strict convexity and symmetry imply $C>0$.

For clarity, write $s=s_\varepsilon$ and denote the six distinct
cluster costs by
\begin{align*}
 a&=G(c,w;\varepsilon,\delta),
 &b&=G(c,w;\varepsilon,1-\delta),\\
 r&=2w\{\phi(\delta)-\phi(1/2)\},
 &q&=2c\{\phi(\varepsilon)-\phi(1/2)\},\\
 t&=r+G(c,2w;\varepsilon,1/2),
 &u&=q+G(2c,w;1/2,\delta).
\end{align*}
Here $a,b$ are the near and far rare-anchor pair costs, $r,q$ are
the rare-pair and anchor-pair costs, and $t,u$ are the triple costs
with respectively one and two anchors. These expressions follow by
first merging the pair inside a triple and then merging its
centroid with the remaining point; expanding both sides proves
this aggregation identity exactly.

By \eqref{eq:gen-remainder-choice},
\[
 a=w\{\phi(\delta)-\phi(\varepsilon)
                  +s(\delta-\varepsilon)\}+o(w)
       \sim w\delta s,
 \qquad b\sim ws.
\]
The first equivalence uses boundedness of $\phi$ on $[0,1]$,
$\varepsilon s\to0$, and $\delta s\to\infty$; the second uses
the same bounds and $\delta\to0$. Continuity at the endpoints
gives $r\sim2wC$ and $q\to C$. Also,
\begin{align*}
 t
 &=2w\{\phi(\delta)-\phi(\varepsilon)\}
          +ws(1-2\varepsilon)+o(w)
   \sim ws.
\end{align*}
Finally, symmetry and differentiability imply $\phi'(1/2)=0$.
Applying \eqref{eq:gen-small-weight} at $1/2$ and using convexity
yields
\[
 0\le u-q\le wd_\phi(\delta\Vert1/2)
       =w\{\phi(\delta)-\phi(1/2)\}\le wC.
\]
Consequently $u\to C$.

We now check all possible partitions rather than assuming that a
particular merger algorithm is optimal. The seven two-cell costs
are $2a,2b,q+r,t,t,u,u$, and the six three-cell costs are
$a,a,b,b,r,q$. The limits just proved show that, for sufficiently
small $\varepsilon$,
\[
 D_{\phi,2}^*=2a,\qquad D_{\phi,3}^*=r.
\]
They also show $b>a>r$, $q>a$, $t>2a$, and $t<q+r$.
If a hierarchy merges the rare pair at level three, its possible
level-two costs are $t,t,q+r$, so its best factor is $t/(2a)$.
If it merges a near rare-anchor pair, it can attain the optimal
level-two partition and has best factor $a/r$. Every other
level-three pair has cost at least $a$, hence cannot improve on
the latter factor. We have therefore proved the exact eventual
formula
\[
 \rho_\phi=\min\{a/r,t/(2a)\}.
\]
Its two entries satisfy
\[
 \frac a r\sim\frac{\delta s}{2C}
            =\frac{\sqrt s}{2\sqrt C},
 \qquad
 \frac{t}{2a}\sim\frac1{2\delta}
            =\frac{\sqrt s}{2\sqrt C}.
\]
This proves \eqref{eq:gen-rate}.
\end{proof}

The choice of $w$ matters. For a general differentiable generator,
setting $w=\varepsilon^2$ without the further reduction in
\eqref{eq:gen-remainder-choice} need not justify the Taylor
remainder. The logarithmic constructions admit explicit estimates
and therefore do not require this adaptive choice. Moreover,
Theorem~\ref{thm:gen-endpoint} is a sufficient condition for an
unbounded hierarchy factor; it is not a characterization of all
generators with this property.

\begin{proposition}[A comparison on intervals of bounded curvature]
\label{prop:gen-curvature}
Let all posteriors lie in a fixed compact interval
$J\subset(0,1)$, and suppose $\phi$ is twice continuously
differentiable there with
$0<m\le\phi''(x)\le M<\infty$. Define weighted squared-error
distortion by
\[
 E(\mathcal P)=\sum_{C\in\mathcal P}
                  \sum_{x\in C}p_x(\eta_x-\bar\eta_C)^2.
\]
Then every partition satisfies
\begin{equation}
 \frac m2 E(\mathcal P)\le D_\phi(\mathcal P)
                         \le\frac M2 E(\mathcal P).
 \label{eq:gen-curvature-comparison}
\end{equation}
For any collection of cardinalities whose optimal distortions
are positive, the minimum simultaneous hierarchy factors obey
\[
 \frac mM\rho_E\le\rho_\phi\le\frac Mm\rho_E.
\]
\end{proposition}

\begin{proof}
The integral form of Taylor's theorem is
\[
 d_\phi(z\Vert x)
 =(z-x)^2\int_0^1(1-t)\phi''(x+t(z-x))\,dt.
\]
Both $x,z$ and their intervening segment belong to $J$, so this
expression lies between $m(z-x)^2/2$ and $M(z-x)^2/2$.
Summing at the cell centroids proves
\eqref{eq:gen-curvature-comparison}. Minimizing its two bounds
over all $k$-cell partitions gives the corresponding bounds on
optimal distortions. Dividing the partition bounds by these
optimal bounds, taking the maximum over cardinalities, and then
the minimum over hierarchies proves the last assertion.
\end{proof}

This comparison transfers hierarchy guarantees between the two
losses up to the fixed condition ratio $M/m$. It does not supply
an independent numerical constant for squared-error hierarchies.
